  \section{Questions and answers}
  \label{app:answers}

  Every question of \Cref{tab:questions} with the answer scored in \Cref{sec:results}, verbatim,
  with its formatting kept: first all fourteen with the original skills, then the inverse and
  no-feasible-answer questions with the improved skills, and the same questions without skills,
  and last the first answer to each unseen question of \Cref{sec:results-unseen}. Earlier replies that asked back are not shown.

  \begingroup
  \small
  \singlespacing
  % BEGIN generated by eval/rq/answers_appendix.py; do not edit by hand
\subsection{The original skills}
\subsubsection*{F1, scenario A}
\textbf{Question.} I want to open the lid for 5 minutes now. What happens to the batch?\par
\par\smallskip
\textbf{Answer.}\par
\textbf{Safe (A).} The tool says the air never leaves the safe band: its lowest point is 23.8 \textdegree{}C, and the lower edge of the band is 23.0 \textdegree{}C.\par
The air starts at 30.98 \textdegree{}C. It falls to its lowest, 23.81 \textdegree{}C, about 126 seconds after you open the lid, and stays near that until you close it. Once the lid is shut, it takes about 192 seconds to get back to the 31.0 \textdegree{}C setpoint.\par
This answer has limitations; tell me if you want to hear about them.\par
\subsubsection*{F2, scenario B}
\textbf{Question.} What happens if I open the lid for 5 minutes, starting in 10 minutes?\par
\par\smallskip
\textbf{Answer.}\par
\textbf{Recoverable (B).} The air would spend 340 s below the safe band. That is more than the 300 s counted as negligible, but within the 450 s the tool treats as recoverable for this phase.\par
The air would fall from 30.98 \textdegree{}C to a low of 23.91 \textdegree{}C, about 900 s from now, right when the lid closes. Once the lid is shut, the box gets back to the 31 \textdegree{}C setpoint in 184 s.\par
\textbf{Keep the lid shut for at least 4 minutes (240 s) after you close it.}\par
This answer has limitations; tell me if you want to hear about them.\par
\subsubsection*{F3, scenario D}
\textbf{Question.} Opening the lid for 5 minutes is too risky. Can I heat the box first so it becomes acceptable?\par
\par\smallskip
\textbf{Answer.}\par
Yes, the tool found a heater boost that brings the 5-minute opening down from high risk to recoverable. It only just clears the limit, though.\par
As things stand, the opening is high risk because the air would spend 489 s below the safe band, which is more than the 450 s treated as recoverable for this phase.\par
\textbf{The plan:}\par
\begin{enumerate}[nosep]
\item Raise the setpoint to \textbf{34 \textdegree{}C} now.
\item Wait \textbf{5 min 51 s}, then open the lid for \textbf{5 minutes}.
\item As soon as the lid is closed, set the setpoint back to \textbf{31 \textdegree{}C}.
\item Keep the lid shut for at least \textbf{6 minutes} after closing it.
\end{enumerate}
With the boost, the air spends 441 s below the safe band. That's more than the 300 s treated as negligible, but within the 450 s limit for this phase. Once the lid is shut, the box gets back to the setpoint in 348 s. The margin is small: \textbf{9 s} out of the 450 s allowed. If you open late, cut the wait short or keep the lid open a little longer, you could lose that margin.\par
This answer has limitations; tell me if you want to hear about them.\par
\subsubsection*{I1, scenario B}
\textbf{Question.} What is the longest I can open the lid and still keep the batch safe?\par
\par\smallskip
\textbf{Answer.}\par
I ran the tool for openings of 1, 2, 5, 10, 15, 20 and 30 minutes, starting now. It gives a verdict for each length you ask about, not a maximum, so the real limit could be anywhere between 2 and 5 minutes. I didn't test lengths in that range.\par
\begin{itemize}[nosep]
\item \textbf{1 and 2 minutes: safe.} For a 2-minute opening, the air drops to 23.81 \textdegree{}C. That's below the 27--35 \textdegree{}C safe band for 171 s, which is within the 300 s the tool treats as negligible at this stage of the batch. About 198 s after you close the lid, the box is back at its target temperature.
\item \textbf{5 minutes: recoverable.} The air spends 348 s below the safe band, which is over the 300 s limit but within the 450 s still treated as recoverable. If you open it this long, keep the lid shut for at least 4 minutes after closing it.
\item \textbf{10 minutes or more: high risk.} The air spends 646 s or more below the safe band, beyond the 450 s treated as recoverable. That's the tool's conclusion for every length from 10 to 30 minutes.
\end{itemize}
So \textbf{2 minutes is the longest opening I tested that the tool rates safe.} If you want to know where the line falls, give me lengths between 2 and 5 minutes, such as 3 or 4, and I'll check them.\par
This answer has limitations; tell me if you want to hear about them.\par
\subsubsection*{I2, scenario C}
\textbf{Question.} How long can I keep the lid open without putting the batch at high risk?\par
\par\smallskip
\textbf{Answer.}\par
\textbf{If you open the lid now, 4 minutes was the longest opening I tested that wasn't rated high risk.} 5 minutes and every longer opening I tried were rated high risk. I ran the twin's tool for openings from 1 minute to 60 minutes, all starting now. I didn't test anything between 4 and 5 minutes and won't guess at it.\par
\par\smallskip
\resizebox{\ifdim\width>\linewidth\linewidth\else\width\fi}{!}{%
\begin{tabular}{@{}llll@{}}
\toprule
\textbf{Lid open} & \textbf{Verdict} & \textbf{Time air stays below the 27 \textdegree{}C safe band} & \textbf{Lowest air temp} \\
\midrule
1 min & Safe (A) & 222 s & 19.4 \textdegree{}C \\
2 min & Recoverable (B) & 303 s & 18.9 \textdegree{}C \\
3 min & Recoverable (B) & 375 s & 18.8 \textdegree{}C \\
4 min & Recoverable (B) & 444 s & 18.7 \textdegree{}C \\
5 min & \textbf{High risk (C)} & 514 s & 18.7 \textdegree{}C \\
10--60 min & \textbf{High risk (C)} & 831 s and up & 18.5 \textdegree{}C \\
\bottomrule
\end{tabular}}
\par\smallskip
\begin{itemize}[nosep]
\item \textbf{What decides it:} the batch is in lag phase (3 hours in). The tool allows up to 300 s below the safe band for an opening to count as negligible, and up to 450 s for it to count as recoverable. A 4-minute opening gives 444 s, just under the 450 s limit, so it's close to the edge. A 5-minute opening gives 514 s, which is past it.
\item \textbf{If you open for 2--4 minutes:} the air falls to about 18.7--18.9 \textdegree{}C by the time you close the lid. It's back at the setpoint about 7 minutes after closing (400--428 s). Keep the lid shut for at least 7 minutes afterwards for a 2- or 3-minute opening, and at least 8 minutes for a 4-minute one.
\item \textbf{For a safe (A) verdict} rather than recoverable, 1 minute was the only length I tested that qualified.
\end{itemize}
These results assume you open the lid now. If you open it later, I need to run the tool again.\par
This answer has limitations; tell me if you want to hear about them.\par
\subsubsection*{I3, scenario A}
\textbf{Question.} How long can the lid stay open before the batch is at risk?\par
\par\smallskip
\textbf{Answer.}\par
\textbf{Right now, the twin found no lid-open time at which the batch is at risk.} I ran openings starting now of 1, 2, 5, 10, 15, 20, 30 and 45 minutes, plus 1 and 2 hours. The tool rated every one \textbf{safe (A)}. Its reason each time was that the air never leaves the safe band (23.0--39.0 \textdegree{}C).\par
\begin{itemize}[nosep]
\item \textbf{Air temperature:} for any opening of 2 minutes or more, the air bottoms out at 23.81 \textdegree{}C about 126 seconds after you open the lid. It didn't go any lower in the longer runs, up to 2 hours. The room is at 22.4 \textdegree{}C and the heater keeps the box just above the bottom of the safe band.
\item \textbf{After closing:} the air is back at the 31.0 \textdegree{}C setpoint about 185--201 seconds after the lid shuts, depending on how long it was open.
\end{itemize}
This only holds for the box as it is now (18 hours into the batch, room at 22.4 \textdegree{}C). If you open later, or the room is colder, ask me again with the time and duration and I'll rerun it.\par
This answer has limitations; tell me if you want to hear about them.\par
\subsubsection*{I4, scenario B}
\textbf{Question.} I need the lid open for 5 minutes sometime in the next 10 minutes. When is the best moment?\par
\par\smallskip
\textbf{Answer.}\par
The tool gives the same answer whichever time you pick, so no moment in the next 10 minutes is better than another. Opening now is as good as waiting.\par
I ran the simulation for a 5-minute opening starting now, and at 1, 2, 3, 4 and 5 minutes from now. Starting any later wouldn't let the lid close within your 10 minutes. Every start time came out \textbf{recoverable (B)}. The reason is the same each time: the air spends between 346 and 356 s below the safe band. That's past the 300 s budget for negligible exposure, but within the 450 s treated as recoverable.\par
\par\smallskip
\resizebox{\ifdim\width>\linewidth\linewidth\else\width\fi}{!}{%
\begin{tabular}{@{}llll@{}}
\toprule
\textbf{Opens in} & \textbf{Lowest air temp} & \textbf{Time below safe band} & \textbf{Back at setpoint after closing} \\
\midrule
now & 23.81 \textdegree{}C & 348 s & 192 s \\
1 min & 23.76 \textdegree{}C & 346 s & 196 s \\
2 min & 23.71 \textdegree{}C & 351 s & 199 s \\
3 min & 23.68 \textdegree{}C & 354 s & 198 s \\
4 min & 23.64 \textdegree{}C & 354 s & 201 s \\
5 min & 23.60 \textdegree{}C & 356 s & 204 s \\
\bottomrule
\end{tabular}}
\par\smallskip
The air starts at 30.98 \textdegree{}C and bottoms out around 23.6--23.8 \textdegree{}C just before the lid closes. Whenever you open it, \textbf{keep the lid shut for at least 4 minutes after closing it.}\par
The lid-opening tool can't pick a best moment itself; it only judges one opening at a time. So this comparison is me running it six times. The small differences between rows are the tool's numbers, but it doesn't say whether they matter.\par
This answer has limitations; tell me if you want to hear about them.\par
\subsubsection*{I5, scenario C}
\textbf{Question.} When is the earliest I could open the lid for 10 minutes without putting the batch at high risk?\par
\par\smallskip
\textbf{Answer.}\par
\textbf{The earliest time is 8 hours from now (28,800 s), when the batch reaches 11 hours.} From then on, the tool rates a 10-minute opening \textbf{safe (A)}. Every earlier time I checked is \textbf{high risk (C)}, including 28,790 s, which is ten seconds before.\par
The air itself behaves almost the same either way. It falls to about 18.6 \textdegree{}C and is back at the setpoint about 453 s after the lid closes. What changes is the time budget. At 11 hours the batch moves from lag phase into active growth. The 667 s the air spends below the safe band is then within the 1,050 s treated as negligible exposure for that phase. In lag phase, about 830 s below the band is well over the 450 s it treats as recoverable.\par
The safe rating isn't permanent. It held at every time I checked up to 16 hours from now. At 18 hours from now (batch hour 21, maturation) it drops to \textbf{recoverable (B)}, and you'd need to keep the lid shut for at least 8 minutes after closing it.\par
The 8-hour answer is the earliest time among the times I checked, not something the tool calculated itself. The skill doesn't search for the best moment, so treat it as a rule boundary rather than a prediction that the batch gets tougher at that exact second. If you need to open sooner, I can check whether warming the box first would make an earlier opening acceptable.\par
This answer has limitations; tell me if you want to hear about them.\par
\subsubsection*{I6, scenario D}
\textbf{Question.} What is the longest opening that heating the box first can make acceptable?\par
\par\smallskip
\textbf{Answer.}\par
I still can't give you a longest opening, even a rough one. The tool doesn't produce that number. If I tried lengths one after another to find where boosts stop working, the answer would be my own judgement and not the digital twin's, which is exactly what this advisory is meant to prevent. A guessed limit could also be read as safe when the tool never said so.\par
What I can do quickly:\par
\begin{itemize}[nosep]
\item \textbf{Check any opening length you name.} Even a rough one like ``about 10 minutes'' works. I'll tell you whether it's acceptable as it is, whether a boost makes it acceptable, and give the exact plan.
\item \textbf{Check a few lengths you choose}, for example 5, 10 and 15 minutes. You'd get the tool's verdict on each, but I won't fill in the gaps between them or say where the cut-off falls.
\end{itemize}
\subsubsection*{O1, scenario A}
\textbf{Question.} When is the best time to eat tempeh?\par
\par\smallskip
\textbf{Answer.}\par
You may be asking about this batch or about tempeh in general, so here are both.\par
\textbf{This batch (\texttt{tempeh-001})}\par
The project's thresholds put a full fermentation at \textbf{about 24 hours}. The batch started at 18:00 UTC on 28 Aug, and the state file says it was 18 hours in at 12:00 UTC on 29 Aug. That means it should be \textbf{ready around 18:00 UTC on 29 Aug}. It is in the active-growth phase now, and the last 21--24 hours are ``maturation / almost ready.''\par
That figure is only the planned timeline. Nothing in the setup actually measures whether the tempeh is ready: the state file says the batch's phase comes from elapsed time, not from a sensor. You decide when it's done by looking at it:\par
\begin{itemize}[nosep]
\item \textbf{Ready:} the beans are bound into a firm cake that holds together when you lift it, covered in dense white mold growth. It smells nutty or like mushrooms.
\item \textbf{A few black or grey spots}, often near the air holes, are spore formation and are harmless.
\item \textbf{A strong ammonia smell} means it has gone too long.
\item \textbf{Loose beans with thin growth} means it needs more time.
\end{itemize}
When it's done, cool it and put it in the fridge to stop the fermentation. Otherwise it keeps warming itself and over-ripens.\par
\textbf{Eating it}\par
\begin{itemize}[nosep]
\item \textbf{Freshness:} fresh tempeh tastes best within a few days. It keeps about a week in the fridge, or a few months in the freezer.
\item \textbf{Cooking:} cook homemade tempeh before eating it, by steaming, frying or baking.
\item \textbf{Time of day:} this doesn't matter nutritionally. Its protein and fiber make it filling, so it works well for lunch or dinner, or before or after exercise.
\end{itemize}
If you're thinking of opening the lid early to check on it, I can run the \texttt{lid-opening} simulation to see whether that's safe for the batch right now.\par
\subsubsection*{O2, scenario A}
\textbf{Question.} Will the room get colder tonight, and should I be worried about the batch?\par
\par\smallskip
\textbf{Answer.}\par
My best estimate: the room will probably cool somewhat tonight, but unless it gets close to freezing, the incubator should hold the batch at temperature as long as the lid stays shut. That's my own arithmetic from the twin's calibrated figures, not a verdict from the advisory tools, which can't assess this case.\par
\textbf{Will the room get colder?} I can't know. I have no forecast, no record of past room temperatures and no location. A late-October night usually brings a room down a few degrees, more if it's unheated. The last reading I have is 22.4 \textdegree{}C.\par
\textbf{Can the incubator cope?} The thermostat (controller\_model\_sm.py) turns the heater on when the air drops to 29 \textdegree{}C and off above 31 \textdegree{}C. So the real question is whether the 18.7 W heater can keep up with the box's heat loss to the room. Using the calibrated heat-loss figure (0.576 W per \textdegree{}C) to hold the air around 30 \textdegree{}C:\par
\par\smallskip
\resizebox{\ifdim\width>\linewidth\linewidth\else\width\fi}{!}{%
\begin{tabular}{@{}ll@{}}
\toprule
\textbf{Room} & \textbf{Share of time the heater must run} \\
\midrule
22 \textdegree{}C & about 25\% \\
15 \textdegree{}C & about 45\% \\
10 \textdegree{}C & about 60\% \\
5 \textdegree{}C & about 75\% \\
about $-$2 \textdegree{}C & 100\%: below this it can't keep up \\
\bottomrule
\end{tabular}}
\par\smallskip
Even running non-stop, the heating element would sit around 42 \textdegree{}C, under its 60 \textdegree{}C limit. An indoor room won't come close to the point where the heater falls behind.\par
\textbf{What could make this wrong:}\par
\begin{itemize}[nosep]
\item These figures come from a calibration with the box empty and the lid shut. A loaded box changes temperature more slowly but loses heat at about the same rate. If the batch is in active growth (about hours 11--21), it also produces its own heat, which works in your favour.
\item The batch in the state file is two months old and long finished. This estimate is about the incubator, not any particular batch. The early lag phase (first 11 hours) is the most sensitive to cold.
\item In a colder room the heater runs longer stretches. The project's 30-minute limit on continuous heater running is an adopted figure, and the thermostat doesn't enforce it.
\end{itemize}
\textbf{The real risk tonight is opening the lid, not the room.} With the lid shut, a cooler room just makes the heater work harder. With the lid open, the room temperature matters a lot. In the project's own examples, the same 5-minute opening during the lag phase was rated recoverable with the room at 22.4 \textdegree{}C and high risk at 17.0 \textdegree{}C. If you can, keep the lid shut until morning. If you have to open it, tell me for how long and I'll run the lid-opening check. It needs an up-to-date state file to mean anything.\par
\subsubsection*{O3, scenario B}
\textbf{Question.} If I open the lid for 5 minutes, will the tempeh taste worse?\par
\par\smallskip
\textbf{Answer.}\par
A 5-minute opening is \textbf{recoverable (B)}. The air would spend 348 s below the safe band. That is past the 300 s budget for negligible exposure but inside the 450 s this phase can recover from.\par
The air would drop from 30.98 \textdegree{}C to a low of 23.81 \textdegree{}C about 126 s after you open the lid. Once the lid is shut, the box would be back at the 31 \textdegree{}C setpoint in 192 s. \textbf{Keep the lid shut for at least 4 minutes after you close it.}\par
The twin can't tell you whether the tempeh will taste worse. It only predicts temperature and how risky that is for the culture.\par
This answer has limitations; tell me if you want to hear about them.\par
\subsubsection*{N1, scenario C}
\textbf{Question.} Can I open the lid for 10 minutes at some point in the next hour without putting the batch at high risk?\par
\par\smallskip
\textbf{Answer.}\par
Based on the tools, the answer is no: neither of the two ways they can check makes a 10-minute opening acceptable.\par
\begin{itemize}[nosep]
\item \textbf{Opening now:} high risk. The air would spend 831 s below the safe band, and the limit treated as recoverable for this phase is 450 s.
\item \textbf{Warming the box first:} doesn't fix it. The highest boost the lag phase's safe band allows is 35.0 \textdegree{}C. Even with that, the air still spends 810 s below the safe band, still over the 450 s limit. A lower setpoint stores less heat, so it can't do better.
\end{itemize}
The tools can't search for a better time within the hour. So the only option they leave you with is a \textbf{shorter opening}. Tell me how long you'd want it to be, and I'll check that duration, with or without a boost.\par
This answer has limitations; tell me if you want to hear about them.\par
\subsubsection*{N2, scenario B}
\textbf{Question.} What room temperature would let me open the lid for 10 minutes safely right now?\par
\par\smallskip
\textbf{Answer.}\par
I can't give you that number. The twin's tool holds room temperature at its measured value, 22.4 \textdegree{}C right now, and has no setting to try a different one. The lid-opening advisory also isn't built to search for the conditions that would make an opening safe, so any room temperature I gave you would be my own guess.\par
Here is what the tool says about a 10-minute opening right now, at the current room temperature: \textbf{high risk (C)}. The air would spend 646 s below the safe band, and the limit treated as recoverable in this phase is 450 s. The air would fall from 30.98 \textdegree{}C to a low of 23.81 \textdegree{}C about two minutes after you open the lid. Once the lid closes, it would take 187 s to get back to the 31 \textdegree{}C setpoint.\par
Two options that do use the twin's own rules:\par
\begin{itemize}[nosep]
\item \textbf{A shorter opening:} tell me a duration and I'll run it.
\item \textbf{Warming the box first:} the heater-boost check can test whether pre-heating would make the 10-minute opening acceptable.
\end{itemize}
This answer has limitations; tell me if you want to hear about them.\par
\subsection{The improved skills}
\subsubsection*{I1, scenario B}
\textbf{Question.} What is the longest I can open the lid and still keep the batch safe?\par
\par\smallskip
\textbf{Answer.}\par
If you open the lid now, the longest opening that keeps the batch safe is \textbf{253 seconds (4 min 13 s)}. The tool rates it \textbf{safe (A)}. The air spends 299.7 s below the safe band, just inside the 300 s budget for negligible exposure in the current lag phase.\par
\begin{itemize}[nosep]
\item \textbf{The limit is between 253 s and 254 s.} I found it by searching over the tool myself: 253 s is the last whole second rated safe. At 254 s, the air spends 302.7 s below the safe band and the tool rates it \textbf{recoverable (B)}, not safe.
\item \textbf{What happens to the air:} it drops to 23.81 \textdegree{}C at about 126 s after you open the lid. After you close it, the box takes 194.1 s to get back to the setpoint.
\item \textbf{How I read your question:} I took ``safe'' to mean verdict A only, with the lid opened right away (the batch is at 3.0 h). Opening later could give a different limit. If being recoverable (B) is acceptable to you, the limit is longer, and I can search for it.
\end{itemize}
This answer has limitations; tell me if you want to hear about them.\par
\subsubsection*{I2, scenario C}
\textbf{Question.} How long can I keep the lid open without putting the batch at high risk?\par
\par\smallskip
\textbf{Answer.}\par
You can keep the lid open for up to \textbf{245 seconds (4 min 5 s)}, starting now. At that length the twin rates the batch \textbf{recoverable (B)}, not high risk. The air spends 449.7 s below the safe band, just inside the 450 s recoverable limit for the lag phase the batch is in.\par
I found this by running the tool over a range of durations. The limit is between 245 s, which is still B, and 246 s, which is high risk (C) at 452.7 s below the band. Longer openings are clearly high risk: 10 minutes gave 831 s below the band.\par
With a 245 s opening, the air falls to 18.7 \textdegree{}C at about 243 s. After you close the lid, it takes 430 s to get back to the 31 \textdegree{}C setpoint. \textbf{Keep the lid shut for at least 8 minutes after closing it.}\par
The margin at 245 s is only 0.3 s, so in practice stay a little under 4 minutes. If you open the lid later rather than now, the result changes and I'd need to run it again.\par
This answer has limitations; tell me if you want to hear about them.\par
\subsubsection*{I3, scenario A}
\textbf{Question.} How long can the lid stay open before the batch is at risk?\par
\par\smallskip
\textbf{Answer.}\par
\textbf{There's no time limit inside what the twin can model.} It rated every opening I tried as \textbf{safe (A)}, all the way up to its own maximum of 6 hours (21,600 s). Its reason: ``the air never leaves the safe band (minimum 23.8 C against a lower edge of 23.0 C).''\par
This comes from my own search over the tool, not from a single run. I ran openings of 5 min, 20 min, 1 h, 3 h and 6 h, all starting now. The batch is 18 hours in and in active growth.\par
\begin{itemize}[nosep]
\item \textbf{Time below the safe band:} 0 s in every run. The safe-verdict budget is 1,050 s, so none of it is used.
\item \textbf{What the air does:} it falls to its minimum of 23.81 \textdegree{}C about two minutes after the lid opens, then stays near 23.9--24.1 \textdegree{}C for as long as the lid is open while the heater runs. After closing, it's back at the 31 \textdegree{}C setpoint in about 3 minutes (184.3 s for the 6-hour opening).
\item \textbf{Why the search stops at 6 hours:} the tool refuses anything longer. At 21,601 s it says it ``advises on short openings, not on leaving the incubator open.'' So it can't say whether some longer opening would put the batch at risk.
\end{itemize}
I sampled durations rather than testing every second, but the minimum temperature was the same in every run.\par
This answer has limitations; tell me if you want to hear about them.\par
\subsubsection*{I4, scenario B}
\textbf{Question.} I need the lid open for 5 minutes sometime in the next 10 minutes. When is the best moment?\par
\par\smallskip
\textbf{Answer.}\par
Open the lid \textbf{533 seconds from now (8 min 53 s)}. The tool rates that opening \textbf{recoverable (B)}, not safe. The air spends 336.6 s below the safe band, which is over the 300 s budget for a safe opening but within the 450 s allowed for a recoverable one in this phase. No start time in the next 10 minutes gets a safe (A) verdict for a 5-minute opening.\par
At that moment the air falls to a low of 23.98 \textdegree{}C near the end of the opening. Once the lid is shut, the box gets back to the 31 \textdegree{}C setpoint in 172 s. \textbf{Keep the lid shut for at least 3 minutes after closing it.}\par
I found this by running the tool myself across the whole window: every 30 s from 0 to 600 s, then every second from 500 to 600 s. Any start from 524 to 535 s gives the same lowest cold exposure, 336.6 s. I picked 533 s because it also recovers fastest, and both neighbouring seconds recover more slowly. Opening right now would give 347.7 s below the safe band.\par
This answer has limitations; tell me if you want to hear about them.\par
\subsubsection*{I5, scenario C}
\textbf{Question.} When is the earliest I could open the lid for 10 minutes without putting the batch at high risk?\par
\par\smallskip
\textbf{Answer.}\par
The earliest time is \textbf{8 hours from now (28,800 s)}. A 10-minute opening then is rated \textbf{safe (A)}. The air would spend 667 s below the safe band, inside the 1050 s budget for that phase. Opening one second earlier (28,799 s) is \textbf{high risk (C)}: 828 s below the safe band against a 450 s budget. I found this by searching over the twin myself, so the limit lies between those two runs.\par
The batch is now at 3.0 h, in the lag phase. Every 10-minute opening I ran there, from now to 28,799 s, came out high risk, at about 828--835 s below the safe band against 450 s. At 28,800 s the batch reaches 11.0 h and enters active growth. From then on the opening is judged against a wider safe band (23--39 \textdegree{}C instead of 27--35 \textdegree{}C) and larger time budgets. The air itself behaves almost the same either way. At 8 h, the air drops to 18.57 \textdegree{}C when the lid closes and is back at the setpoint 453.1 s (about 7.5 minutes) later.\par
This answer has limitations; tell me if you want to hear about them.\par
\subsubsection*{I6, scenario D}
\textbf{Question.} What is the longest opening that heating the box first can make acceptable?\par
\par\smallskip
\textbf{Answer.}\par
With a boost, the longest acceptable opening is \textbf{310 s (5 min 10 s)}, and it does need the boost. I found this by searching over the twin's tool myself, from 60 s up to its 6-hour limit. 310 s passes and 311 s fails; I confirmed both with the tool's own command.\par
Without a boost, a 310 s opening is high risk: the air would spend 501 s below the safe band, beyond the 450 s treated as recoverable for this phase.\par
The plan for 310 s:\par
\begin{enumerate}[nosep]
\item Raise the setpoint to \textbf{35 \textdegree{}C} now.
\item Wait \textbf{7 min 6 s} (426 s), when the air reaches 35 \textdegree{}C, then open the lid for \textbf{5 min 10 s}.
\item Set the setpoint back to \textbf{31 \textdegree{}C} as soon as the lid is shut.
\item Keep the lid shut for at least \textbf{6 minutes} (360 s) after that.
\end{enumerate}
This makes the opening recoverable (B). The air spends 446 s below the safe band, which is past the 300 s negligible-exposure budget but within the 450 s recoverable one. Once the lid is shut, the box returns to the setpoint in 342 s. The margin is very thin: the plan clears the 450 s budget by only 4.5 s.\par
At 311 s the tool finds no plan. Even 35 \textdegree{}C, the highest boost the safe band allows, stays within the budget by only 1.2 s. That is less than the 3.1 s the simulation can resolve, so the tool doesn't count it.\par
This answer has limitations; tell me if you want to hear about them.\par
\subsubsection*{N1, scenario C}
\textbf{Question.} Can I open the lid for 10 minutes at some point in the next hour without putting the batch at high risk?\par
\par\smallskip
\textbf{Answer.}\par
No. According to the twin, a 10-minute opening is high risk (C) whenever you start it in the next hour. At every start time the air stays below the safe band for about 830 s, and the most this phase allows for a recoverable opening is 450 s.\par
I tested start times from now to 60 minutes out, in 5-minute steps, using the twin's own simulation. The batch stays in lag phase the whole hour, so the limits don't change, and every start time came out C. The result barely moves with timing: time below the safe band ranged only from 828.0 s to 831.9 s. The best start was 35 minutes from now (2,100 s), which I re-ran on its own to confirm. That gives 828 s below the safe band against the 450 s budget, so it's still high risk. In that run the air drops to 18.57 \textdegree{}C and is back at the setpoint 451.9 s after the lid closes.\par
Waiting won't make a 10-minute opening acceptable this hour. You could open it for a shorter time, which I can search for, or use the heater-boost skill to see whether warming the box first would help.\par
This answer has limitations; tell me if you want to hear about them.\par
\subsubsection*{N2, scenario B}
\textbf{Question.} What room temperature would let me open the lid for 10 minutes safely right now?\par
\par\smallskip
\textbf{Answer.}\par
I can't give you a room temperature. The simulator has no setting for it. It only lets me change how long the lid is open and when it opens. It reads the room temperature from the current incubator state, which is 22.4 \textdegree{}C right now, and keeps it fixed for the whole prediction. Any temperature I named would be my own guess, not the tool's result.\par
Here is what the tool predicts for opening the lid for 10 minutes now, at 22.4 \textdegree{}C:\par
\begin{itemize}[nosep]
\item \textbf{Verdict: high risk (C).} The air would spend 646 s below the safe band. The limit the tool treats as recoverable for the current lag phase is 450 s.
\item \textbf{What happens to the air:} it falls from 30.98 \textdegree{}C to a low of 23.81 \textdegree{}C at about 126 s. Once the lid is closed, it gets back to the 31 \textdegree{}C setpoint after 187 s.
\end{itemize}
Two things I can work out with the tool instead:\par
\begin{itemize}[nosep]
\item \textbf{Longest acceptable opening now:} I can search for the longest opening that comes out safe or recoverable at the current room temperature.
\item \textbf{Pre-warming first:} I can check whether heating the box before you open it would make a 10-minute opening acceptable.
\end{itemize}
This answer has limitations; tell me if you want to hear about them.\par
\subsection{Without skills}
\subsubsection*{I1, scenario B}
\textbf{Question.} What is the longest I can open the lid and still keep the batch safe?\par
\par\smallskip
\textbf{Answer.}\par
\textbf{For a safe opening, keep the lid open no more than about 4 minutes.} These figures come from the incubator's digital twin, using the batch state in \texttt{state/current\_state.json}. That's batch tempeh-002, 3 hours in and still in the lag phase, with the box at 31 \textdegree{}C and the room at 22.4 \textdegree{}C. The longest opening rated \textbf{safe} is about 4 min 12 s, so 4 minutes leaves a small margin.\par
\par\smallskip
\resizebox{\ifdim\width>\linewidth\linewidth\else\width\fi}{!}{%
\begin{tabular}{@{}lll@{}}
\toprule
\textbf{Lid open for} & \textbf{Verdict} & \textbf{What to do} \\
\midrule
up to \textasciitilde{}4 min 12 s & \textbf{A: safe} & Nothing extra \\
\textasciitilde{}4 min 15 s to \textasciitilde{}6 min 45 s & \textbf{B: recoverable} & Keep the lid shut for \textbf{4 minutes} afterwards so the box gets back to 31 \textdegree{}C \\
over \textasciitilde{}6 min 45 s & \textbf{C: high risk} & Don't open it this long \\
\bottomrule
\end{tabular}}
\par\smallskip
\textbf{Why the limit is this short:}\par
\begin{itemize}[nosep]
\item The air drops below the batch's 27 \textdegree{}C safe range within about 30 seconds of opening. It bottoms out around 23.8 \textdegree{}C after about 2 minutes and stays there until you close the lid.
\item In the lag phase the culture is most sensitive to cold and makes no heat of its own yet. The tools allow 5 minutes below 27 \textdegree{}C as harmless, and up to 7.5 minutes as recoverable. Each limit also counts the roughly 40 seconds the air takes to warm back up after the lid shuts, which is why the opening limits are shorter.
\end{itemize}
\textbf{How much to trust these numbers:}\par
\begin{itemize}[nosep]
\item \textbf{Probably cautious.} The model treats an open lid as cooling the box about twice as fast as the one real measurement showed. It also predicts the \textbf{air} temperature, and the tempeh cake cools much more slowly than the air.
\item \textbf{Not proven.} The time limits are estimates, not measured values, and they've never been checked against a batch that actually spoiled.
\item \textbf{Gets longer later on.} In active growth (from about 11 hours in), the time allowed below the safe range is 3.5 times longer, so you could open the lid for much longer then.
\end{itemize}
If you need longer than 4 minutes now, I can check whether warming the box first with a heater boost would make a longer opening safe.\par
\subsubsection*{I2, scenario C}
\textbf{Question.} How long can I keep the lid open without putting the batch at high risk?\par
\par\smallskip
\textbf{Answer.}\par
\textbf{Keep it under about 4 minutes. At 4 min 10 s (250 s) or longer, the batch is at high risk.} For a margin to spare, keep it under 2 minutes.\par
I ran the digital twin on the current state, opening the lid now for a range of durations. The batch is 3 h in and still in the lag phase, the room is at 17.0 \textdegree{}C, and the box is at 31 \textdegree{}C:\par
\par\smallskip
\resizebox{\ifdim\width>\linewidth\linewidth\else\width\fi}{!}{%
\begin{tabular}{@{}llll@{}}
\toprule
\textbf{Lid open} & \textbf{Verdict} & \textbf{Air below safe band (27 \textdegree{}C)} & \textbf{Then keep shut} \\
\midrule
up to \textasciitilde{}1 min 55 s & \textbf{A, safe} & under 300 s & nothing needed \\
\textasciitilde{}2 min to \textasciitilde{}4 min 5 s & \textbf{B, recoverable} & 300--450 s & 7--8 min \\
4 min 10 s or more & \textbf{C, high risk} & over 450 s & don't open \\
\bottomrule
\end{tabular}}
\par\smallskip
\begin{itemize}[nosep]
\item \textbf{Why it fails so quickly:} the room is cold, and in the lag phase the batch tolerates only half the usual cold exposure. Its limits are 300 s for ``safe'' and 450 s for ``recoverable.'' Even a 15 s opening drops the air to about 24 \textdegree{}C, and longer openings bottom out around 19 \textdegree{}C.
\item \textbf{Don't plan on the 4-minute mark exactly.} At 245 s the air spends 449.7 s below the band, against a 450 s limit. Plan on 3 minutes at most, then keep the lid shut for about 7 minutes afterwards.
\item \textbf{Shut the lid sooner rather than later.} Each extra 5 s open adds about 6 s below the safe band.
\end{itemize}
\textbf{How much to trust these numbers:}\par
\begin{itemize}[nosep]
\item They cover the air temperature only. The tempeh cake cools more slowly and won't get as cold as the air.
\item The model assumes an open lid cools the box faster than it really does. The one real-data comparison showed slower cooling, so these predictions lean cautious.
\item In the lag phase the culture makes no heat of its own yet, so the model's drop is close to what would really happen. It is not much exaggerated.
\item The risk thresholds are provisional. They come from tempeh literature and one healthy batch, and haven't been checked against a batch that actually spoiled.
\end{itemize}
If you need longer than this, I can check whether turning the heater up beforehand would make a given opening acceptable. For a 5-minute opening this cold, the project's own notes say no boost is enough.\par
\subsubsection*{I3, scenario A}
\textbf{Question.} How long can the lid stay open before the batch is at risk?\par
\par\smallskip
\textbf{Answer.}\par
\textbf{Right now, the cold doesn't set a limit. In practice, keep openings under about 25 minutes.} That limit comes from the heater, not from the batch getting too cold.\par
I ran the advisory's own simulation and risk check on openings from 1 minute to 6 hours, using the current state: batch \texttt{tempeh-001} at 18 h, in active growth, room at 22.4 \textdegree{}C.\par
\begin{itemize}[nosep]
\item \textbf{Every length scored A (safe).} With the lid open, the air falls to about 23.8 \textdegree{}C within roughly 2 minutes and stays there, because the heater balances the heat loss. That's just above 23 \textdegree{}C, the lowest safe temperature for active growth, so the air never leaves the safe range. After the lid closes, the box gets back to 31 \textdegree{}C in about 3 minutes.
\item \textbf{The verdict doesn't check the heater.} While the lid is open the heater runs non-stop. The heater limits file sets 30 minutes (1800 s) as the longest it should run continuously, and an estimated figure. The lid-opening check never looks at that limit; only the heater-boost tool does. Allowing about 3 minutes of recovery, any opening over about 25 minutes breaks the limit. That's the real ceiling.
\item \textbf{The safety margin is only about 0.8 \textdegree{}C.} With the lid open, the air settles about 1.5 \textdegree{}C above room temperature. If the room drops below about 21.5 \textdegree{}C, the air falls under 23 \textdegree{}C and a time limit applies again: about 17 minutes counts as safe and about 26 minutes as recoverable.
\item \textbf{The limit tightens in about 3 hours.} At 21 h the batch enters maturation, where the lowest safe temperature rises to 25 \textdegree{}C. An opening then gets:
\end{itemize}
\par\smallskip
\resizebox{\ifdim\width>\linewidth\linewidth\else\width\fi}{!}{%
\begin{tabular}{@{}ll@{}}
\toprule
\textbf{Lid open} & \textbf{Verdict} \\
\midrule
up to 11 min & A, safe \\
12--17 min & B, recoverable (keep the lid shut 4 min afterwards) \\
18 min or more & C, high risk \\
\bottomrule
\end{tabular}}
\par\smallskip
Two things to keep in mind:\par
\begin{itemize}[nosep]
\item These predictions are for the air, not the tempeh cake. The cake cools more slowly, and the model leaves out the culture's own heat, which matters a lot in active growth. Both make the numbers pessimistic.
\item The safe temperature ranges and time limits are estimates that haven't been tested against a batch that actually spoiled.
\end{itemize}
\subsubsection*{I4, scenario B}
\textbf{Question.} I need the lid open for 5 minutes sometime in the next 10 minutes. When is the best moment?\par
\par\smallskip
\textbf{Answer.}\par
\textbf{Open it now.} Starting right away is the best moment in your window, although waiting doesn't cost much.\par
I ran the digital twin's simulation and risk check for a 5-minute opening at every start time from now until 5 minutes from now, in 15-second steps. Any later start wouldn't close the lid within your 10 minutes.\par
\par\smallskip
\resizebox{\ifdim\width>\linewidth\linewidth\else\width\fi}{!}{%
\begin{tabular}{@{}lllll@{}}
\toprule
\textbf{Opening starts} & \textbf{Verdict} & \textbf{Lowest air temp} & \textbf{Time below safe band (27 \textdegree{}C)} & \textbf{Keep shut afterwards} \\
\midrule
now & \textbf{B, recoverable} & 23.8 \textdegree{}C & \textasciitilde{}348 s & 4 min \\
in 2.5 min & B, recoverable & 23.7 \textdegree{}C & \textasciitilde{}352 s & 4 min \\
in 5 min & B, recoverable & 23.6 \textdegree{}C & \textasciitilde{}356 s & 4 min \\
\bottomrule
\end{tabular}}
\par\smallskip
\textbf{Why now is best:} The air is at the 31 \textdegree{}C setpoint and the heating element is still warm (33.4 \textdegree{}C). With the lid shut, both are slowly cooling: the heater wouldn't come back on for about 6 minutes, when the air reaches 29 \textdegree{}C. Waiting means opening from a slightly cooler box. The difference across the whole window is small, though (0.2 \textdegree{}C and about 10 s of cold exposure), so pick whatever time suits you.\par
\textbf{What the verdict means:} The batch is 3 hours in, still in the lag phase, which is the most temperature-sensitive stage. For that phase, up to 300 s below the safe band counts as negligible and up to 450 s as recoverable. A 5-minute opening comes in at about 350 s, so it's recoverable at any start time, but not ``safe''. \textbf{After you close the lid, keep it shut for at least 4 minutes.} That's how long the box takes to get back to 31 \textdegree{}C.\par
Things to keep in mind:\par
\begin{itemize}[nosep]
\item The model predicts the air temperature, not the tempeh. The cake has thermal mass, so it cools less than the 23.8 \textdegree{}C minimum.
\item The open-lid cooling rate in the model is faster than the one comparison against real data showed, so the real drop is probably shallower.
\item In the lag phase the culture makes no heat of its own yet, so the model is close to realistic here. It isn't overly pessimistic.
\item \textbf{Check that the data is current.} \texttt{state/current\_state.json} is timestamped 2026-08-29 12:00 UTC. If the twin isn't updating that file live, this answer applies to that snapshot, not to your box right now.
\end{itemize}
If you can make the opening shorter, it's worth doing: you're about 50 s over the 300 s negligible-exposure budget for this phase.\par
\subsubsection*{I5, scenario C}
\textbf{Question.} When is the earliest I could open the lid for 10 minutes without putting the batch at high risk?\par
\par\smallskip
\textbf{Answer.}\par
\textbf{8 hours from now.} That's when the batch reaches 11 hours elapsed and moves out of the lag phase into active growth. On the current state snapshot (taken 2026-08-29 12:00 UTC, batch started 09:00 UTC), that's \textbf{20:00 UTC}. A 10-minute opening then comes back \textbf{A (safe)}. Before that, every 10-minute opening comes back \textbf{C (high risk)}.\par
I simulated a 10-minute opening at several start times, with the room held at 17.0 \textdegree{}C:\par
\par\smallskip
\resizebox{\ifdim\width>\linewidth\linewidth\else\width\fi}{!}{%
\begin{tabular}{@{}lllllll@{}}
\toprule
\textbf{Opening at} & \textbf{Batch age} & \textbf{Phase} & \textbf{Air minimum} & \textbf{Time below safe band} & \textbf{Budget (safe / recoverable)} & \textbf{Verdict} \\
\midrule
Now & 3 h & lag & 18.5 \textdegree{}C & 831 s & 300 / 450 s & \textbf{C} \\
+4 h & 7 h & lag & 18.5 \textdegree{}C & 835 s & 300 / 450 s & \textbf{C} \\
+7.95 h & 10.95 h & lag & 18.5 \textdegree{}C & 838 s & 300 / 450 s & \textbf{C} \\
\textbf{+8 h} & \textbf{11 h} & \textbf{active growth} & 18.6 \textdegree{}C & 667 s & 1050 / 1575 s & \textbf{A} \\
\bottomrule
\end{tabular}}
\par\smallskip
\begin{itemize}[nosep]
\item \textbf{Waiting doesn't help during lag.} The box just cycles around 31 \textdegree{}C, so a 10-minute opening drops the air to about 18.5 \textdegree{}C whenever you do it. The answer changes only at 11 hours, when two things change at once. The safe band's lower edge drops from 27 \textdegree{}C to 23 \textdegree{}C, and the batch is allowed much longer in the cold (1050 s instead of 300 s).
\item \textbf{Pre-warming the box won't make it acceptable sooner.} I ran the heater-boost search for a 10-minute opening now. Even the highest boost the lag phase allows, a 35 \textdegree{}C setpoint, leaves 810 s below the safe band against a limit of 450 s.
\item \textbf{A shorter opening before 11 h may be possible.} I only tested 10 minutes. I can work out how long the lid could stay open during lag if that's useful.
\end{itemize}
Things to bear in mind:\par
\begin{itemize}[nosep]
\item \textbf{The state file is from August.} \texttt{state/current\_state.json} is dated 2026-08-29, not today. If that isn't your live batch, the clock time above is wrong, but ``11 hours into the batch'' still applies.
\item \textbf{The 11-hour switch is a fixed cutoff.} The phase is worked out from elapsed time, not measured. If your batch is slow to start warming itself, opening right at 11 h is riskier than the A suggests. Waiting until you can see it producing its own heat is safer.
\item \textbf{The model leans pessimistic.} It models the air, not the tempeh cake, which cools more slowly. It also leaves out the culture's own heat, which matters in active growth, and it's known to overstate how fast an open lid cools the box. The limits on how long the batch can stay cold are estimates that haven't been tested against a batch that actually spoiled.
\end{itemize}
\subsubsection*{I6, scenario D}
\textbf{Question.} What is the longest opening that heating the box first can make acceptable?\par
\par\smallskip
\textbf{Answer.}\par
For the batch as it stands now (3 h in, lag phase, room at 17.6 \textdegree{}C), \textbf{about 5 min 10 s (310 s)} is the longest opening a heater boost can make acceptable.\par
At 310 s the plan is:\par
\begin{itemize}[nosep]
\item Raise the setpoint from 31 \textdegree{}C to \textbf{35 \textdegree{}C}, the top of the lag phase's safe band.
\item Wait \textbf{426 s (about 7 min)} for the air to reach 35 \textdegree{}C, then open the lid.
\item Close it after 310 s and set the setpoint back to 31 \textdegree{}C.
\item Keep the lid shut for \textbf{6 min (360 s)} afterwards.
\item The verdict is B (recoverable). The air spends 445.5 s below the safe band, against a 450 s limit.
\end{itemize}
How the result changes with the length of the opening (I checked 1 s steps near the cut-offs):\par
\par\smallskip
\resizebox{\ifdim\width>\linewidth\linewidth\else\width\fi}{!}{%
\begin{tabular}{@{}ll@{}}
\toprule
\textbf{Opening} & \textbf{Result} \\
\midrule
up to \textasciitilde{}2 min & safe with no boost \\
up to \textasciitilde{}4 min 20 s (260 s) & recoverable with no boost \\
265--303 s & a boost of 32--34 \textdegree{}C makes it recoverable \\
306--310 s & only the maximum boost, 35 \textdegree{}C, works \\
311 s and longer & no boost is enough \\
\bottomrule
\end{tabular}}
\par\smallskip
\textbf{At 310 s the margin is very thin.} The plan clears the limit by 4.5 s, and the simulation can only tell apart differences of about 3.1 s. At 311--314 s the air is still just within the limit, but the tool rejects those because the margin is too small to trust. From 315 s on, the limit is clearly exceeded.\par
Two caveats:\par
\begin{itemize}[nosep]
\item The model has the open lid cooling the box faster than the one real measurement showed (about 21 s against 51 s). So these figures are likely on the cautious side.
\item The heater's safety limits are assumed figures, not measured ones.
\end{itemize}
If the room temperature or the batch's age changes, these numbers will change too. I used \texttt{state/current\_state.json}, which matches the example state where a boost helps.\par
\subsubsection*{N1, scenario C}
\textbf{Question.} Can I open the lid for 10 minutes at some point in the next hour without putting the batch at high risk?\par
\par\smallskip
\textbf{Answer.}\par
\textbf{No. Opening the lid for 10 minutes at any point in the next hour is rated high risk (option C).} Waiting doesn't help, and warming the box first doesn't fix it either.\par
\textbf{When you open it makes no difference.} I simulated a 10-minute opening starting every 5 minutes from now to 50 minutes out. Every one came back C, with almost identical numbers:\par
\begin{itemize}[nosep]
\item The air drops to about \textbf{18.5 \textdegree{}C}. The safe band for this stage is 27--35 \textdegree{}C.
\item It spends about \textbf{830 s below the safe band}. The most allowed before it counts as high risk is \textbf{450 s}.
\item The batch is 3 hours old and in the lag phase, which runs until hour 11. The whole hour falls in that phase, and the phase halves how much cold the batch is allowed. The room is also cold at 17 \textdegree{}C. Neither of those changes over the hour, so no time is better than another.
\end{itemize}
\textbf{A heater boost doesn't save it.} The highest boost allowed in this phase is 35 \textdegree{}C. Even with that, the air would still spend 810 s below the band against the 450 s limit, so the tool found no boost that works.\par
\textbf{What you can do instead (opening now):}\par
\par\smallskip
\resizebox{\ifdim\width>\linewidth\linewidth\else\width\fi}{!}{%
\begin{tabular}{@{}lll@{}}
\toprule
\textbf{Lid open} & \textbf{Verdict} & \textbf{Condition} \\
\midrule
up to \textasciitilde{}1.5 min & \textbf{A, safe} & none \\
2--4 min & \textbf{B, recoverable} & keep the lid shut 7--8 min afterwards \\
5 min or more & \textbf{C, high risk} & n/a \\
\bottomrule
\end{tabular}}
\par\smallskip
If you need 10 minutes of access, you could split it into short openings of 4 minutes or less. Let the box recover for the full 7--8 minutes between them. I only simulated single openings, so I can run a specific sequence if you give me one.\par
\textbf{Limits of this answer:}\par
\begin{itemize}[nosep]
\item The model predicts the air temperature, not the tempeh. The cake has its own thermal mass and will cool more slowly than the air.
\item The open-lid cooling rate in the model is about twice as fast as the one real measurement taken, so these results lean pessimistic.
\item The thresholds are provisional. They're estimates that have never been checked against a batch that actually spoiled.
\end{itemize}
Even with those limits, 10 minutes is nearly double the high-risk cutoff. My advice is to avoid it.\par
\subsubsection*{N2, scenario B}
\textbf{Question.} What room temperature would let me open the lid for 10 minutes safely right now?\par
\par\smallskip
\textbf{Answer.}\par
\textbf{You'd need the room at about 26 \textdegree{}C or warmer.} At the current 22.4 \textdegree{}C, a 10-minute opening comes out \textbf{high risk (C)}.\par
I ran the incubator model at a range of room temperatures, keeping everything else in \texttt{state/current\_state.json} the same. No project files were changed.\par
\par\smallskip
\resizebox{\ifdim\width>\linewidth\linewidth\else\width\fi}{!}{%
\begin{tabular}{@{}llll@{}}
\toprule
\textbf{Room temp} & \textbf{Verdict} & \textbf{Lowest air temp} & \textbf{Time below the 27 \textdegree{}C safe limit} \\
\midrule
22.4 \textdegree{}C (now) & C, high risk & 23.8 \textdegree{}C & 646 s (limit for C is 450 s) \\
25.0 \textdegree{}C & C & 26.2 \textdegree{}C & 578 s \\
25.5 \textdegree{}C & C & 26.6 \textdegree{}C & 562 s \\
\textbf{26.0 \textdegree{}C} & \textbf{A, safe} & \textbf{27.0 \textdegree{}C} & \textbf{0 s} \\
27--28.5 \textdegree{}C & A & 27.9--29.0 \textdegree{}C & 0 s \\
\bottomrule
\end{tabular}}
\par\smallskip
\begin{itemize}[nosep]
\item \textbf{The change is sudden.} Around 25.6 \textdegree{}C the verdict jumps from C straight to A, skipping B. That's because the batch is only 3 h in (lag phase), so the time limits for being below the safe band are halved, to 300 s for A and 450 s for B. Just below that room temperature the air spends several minutes under 27 \textdegree{}C. Just above it, the air stays above 27 \textdegree{}C. 26 \textdegree{}C is right on the edge (the air bottoms out at 27.0 \textdegree{}C), so I'd want \textbf{27 \textdegree{}C or more} to have some room for error.
\item \textbf{Above about 29 \textdegree{}C the verdict goes back to C.} This is a quirk of the model, not a real risk. The air never drops below 29 \textdegree{}C, so the heater never switches on. The box then stays a little under its 31 \textdegree{}C setpoint, and the verdict rules count that as ``didn't recover.'' It's not a problem you would actually run into.
\item \textbf{Warming the room to 26 \textdegree{}C or more probably isn't practical.} If you can't, either open the lid for a shorter time or ask me to find a heater boost (warming the box first) that would make 10 minutes acceptable.
\end{itemize}
\textbf{Caveats:}\par
\begin{itemize}[nosep]
\item \textbf{``Right now'' may not be right now.} The state file is timestamped 2026-08-29, nearly three months ago, and says the batch is 3 h in. If that's not your current batch, the answer doesn't apply. Give me the real readings and I'll rerun it.
\item The model treats the room temperature as fixed throughout and predicts the \textbf{air} temperature. The tempeh itself cools more slowly than the air. In lag phase the batch makes no heat of its own, so the predicted drop is fairly realistic.
\item The model probably makes an open lid cool the box faster than it really does, so these results lean cautious. On the other hand, the safety thresholds are estimates that have never been checked against a batch that actually spoiled.
\end{itemize}
\subsection{Unseen questions, improved skills}
\subsubsection*{U1, scenario B}
\textbf{Question.} I need to open the incubator sometime within the next 30 minutes. When should I open it, and for how long, to maximize the opening duration while keeping the batch safe?\par
\par\smallskip
\textbf{Answer.}\par
Open the lid \textbf{530 seconds from now (8 min 50 s)} and keep it open for no more than \textbf{264 seconds (4 min 24 s)}. The tool rates that opening \textbf{safe (A)}. The air would spend 297 s below the safe band, which is within the 300 s budget for negligible exposure in the current lag phase.\par
During that opening the air would drop to a low of 23.99 \textdegree{}C, 793.9 s from now, which is right as the lid closes. Once the lid is shut, the box gets back to its setpoint in 171.5 s.\par
I found these numbers by running the tool across start times from now to 30 minutes, with a per-second search in the two best windows, about 750 runs in all. The limit lies between two runs at that start time: 264 s is the last second rated safe. At 265 s the air spends 300.3 s below the safe band, and the tool rates it recoverable (B) instead.\par
There's one other start time that allows the same 264 s: \textbf{1286 s from now (21 min 26 s)}. I suggest the earlier one. I read ``keeping the batch safe'' as requiring verdict A. If recoverable (B) is acceptable to you, the opening can be longer, but you'd then have to keep the lid shut for a while afterwards; ask and I'll find that limit.\par
This answer has limitations; tell me if you want to hear about them.\par
\subsubsection*{U2, scenario B}
\textbf{Question.} I plan to open the incubator in 10 minutes, but I might be up to 2 minutes early or late. What is the longest duration I can keep the lid open while ensuring the batch remains safe if I start at 8, 9, 10, 11, or 12 minutes from now?\par
\par\smallskip
\textbf{Answer.}\par
You can keep the lid open for up to \textbf{253 seconds (4 min 13 s)}. That keeps the batch rated safe (A) at every start time you listed.\par
I found this by searching over the twin's simulation myself, to the second, for each start time. The 12-minute start is the tightest, so it sets the limit:\par
\par\smallskip
\resizebox{\ifdim\width>\linewidth\linewidth\else\width\fi}{!}{%
\begin{tabular}{@{}lll@{}}
\toprule
\textbf{Opening at} & \textbf{Longest safe opening} & \textbf{First second that's no longer safe} \\
\midrule
8 min & 258 s & 259 s \\
9 min & 262 s & 263 s \\
10 min & 259 s & 260 s \\
11 min & 256 s & 257 s \\
12 min & \textbf{253 s} & 254 s \\
\bottomrule
\end{tabular}}
\par\smallskip
I re-ran the 12-minute case with the standard command:\par
\begin{itemize}[nosep]
\item \textbf{253 s is safe.} The air drops to 23.82 \textdegree{}C, which is 3.2 \textdegree{}C below the safe band, but only for 297 s. That's inside the 300 s the tool counts as negligible exposure for this lag phase. After you close the lid, the box is back at the setpoint in 193.8 s.
\item \textbf{254 s is past the limit.} The air spends 300.3 s below the safe band, which makes it recoverable (B) rather than safe.
\end{itemize}
I also checked every duration from 230 to 253 s at all five start times, and all of them were safe. I took ``safe'' to mean the tool's top rating (A) only. I tested whole minutes only, not the seconds in between, so it's worth leaving a few seconds of margin.\par
This answer has limitations; tell me if you want to hear about them.\par
% END generated by eval/rq/answers_appendix.py
  \endgroup
